\documentclass{situationist-preprint}

\title{Duchamp's Fountain and the Limits of the Institutional Theory}
\msid{AES-01-c}
\msfamily{AES-01}
\msdiscipline{Aesthetics}
\mstype{Consequence}
\mskeywords{Duchamp; Fountain; readymade; institutional theory; critical uptake; replicas}

\begin{document}
\maketitle

\begin{abstract}
Duchamp's \emph{Fountain} is the paradigm case for institutional accounts of art, and also the case most often used against them. We derive what the account of AES-01, on which art is what is presented as art within an artworld subject to the artworld's critical uptake, implies about it. Three consequences follow: \emph{Fountain} is a work of art; its status was not secured by its submission in 1917, which was refused, but by the critical uptake that followed; and the later replicas share its status only through the practice that takes them up as instances of it. Two limits follow as well: the account classifies \emph{Fountain} but does not explain its importance, and it does not settle whether its uptake was warranted.
\end{abstract}

\section{The case}

In April 1917 a porcelain urinal, turned on its back and signed \enquote{R. Mutt 1917}, was submitted to the first exhibition of the Society of Independent Artists in New York, whose rules promised to hang any work by a member who paid the fee. The object was not shown. Shortly afterwards it was photographed by Alfred Stieglitz, and the photograph was reproduced in the second issue of \emph{The Blind Man}, with an unsigned editorial defending it on the ground that the artist had chosen an ordinary article of life and placed it so that its useful significance disappeared under a new title and point of view. The original object was lost. From 1950 onward Duchamp authorized replicas, including an edition produced with Arturo Schwarz in 1964, and these are the objects now held in museums.

The case is central for three reasons. It is an object indiscernible from objects that are not art, which is the phenomenon \citet{Danto1964,Danto1981} took to show that art status cannot be a perceptible property. Its maker did not fashion it, which strains accounts on which art is the product of skill. And it was offered precisely as a test of what an artworld would accept. We take the account of AES-01 as given, with the notions of presentation and critical uptake as defined in Definitions~2.1 and~2.2 of AES-15-a, and derive what it implies.

\section{Consequences}

\begin{proposition}\label{p:art}
\emph{Fountain} is a work of art.
\end{proposition}

\begin{proof}
The object was presented: Duchamp, a participant in the New York artworld, offered it through the Society's exhibition and then through \emph{The Blind Man} to a public as a candidate for the appreciation and criticism the artworld practices. It received critical uptake: it has been interpreted, compared with other works and assessed, adversely as well as favourably, by the critical practice of the artworld continuously since the 1950s and episodically before \citep{deDuve1996}. By the account of AES-01 it is therefore a work of art.
\end{proof}

\begin{proposition}\label{p:refusal}
Submission to the Society did not by itself make \emph{Fountain} a work of art.
\end{proposition}

\begin{proof}
Submission met the presentation condition. But on the account of AES-01, presentation is subject to uptake, and in 1917 the Society's directors declined to exhibit the object, on the ground, as reported, that it was not art. At that point the object had been offered and refused, which is precisely a case in which presentation is met and uptake is not. Its status depended on the engagement that followed: the defense in \emph{The Blind Man}, and the later interpretation that made it a reference point for twentieth-century art.
\end{proof}

Proposition~\ref{p:refusal} corrects a common simplification, on which Duchamp made the urinal art by an act of declaration. The account of AES-01 implies the opposite: no declaration could do so, and the refusal of 1917 is not an embarrassment for the institutional theory but a confirmation that presentation alone does not suffice. The case also shows that the artworld whose uptake matters is not identical with any one institution within it. The Society refused; the practice of criticism, which is wider than the Society and outlasted it, did not.

\begin{proposition}\label{p:replica}
The replicas are works of art only as instances of \emph{Fountain}, or as works in their own right, as the critical practice takes them up.
\end{proposition}

\begin{proof}
The replicas are distinct objects, made decades later and in some cases by craftsmen working from the photograph. They are not the object that was submitted. On the account of AES-01, their status depends on their own presentation and uptake. They were presented, with Duchamp's authorization, as the work, and the practice has engaged them as the work: they are exhibited, catalogued and interpreted as \emph{Fountain}. Their status is therefore derivative from the practice's treating them as instances of a work whose first instance was lost.
\end{proof}

\section{Limits}

\begin{remark}\label{r:import}
The account classifies \emph{Fountain} but does not explain why it matters.
\end{remark}

The account says what makes \emph{Fountain} a work; it says nothing about why it became one of the most discussed works of its century. That explanation is art-historical and critical: it concerns what the object showed about the conventions of exhibition, authorship and craft, and the use later artists made of it. \citet{Binkley1977} argued from such cases that art need not be aesthetic; the institutional account is compatible with that claim but does not establish it, and its silence on importance is a consequence of its being an account of status rather than of value.

\begin{remark}\label{r:warrant}
The account does not settle whether the uptake of \emph{Fountain} was warranted.
\end{remark}

A critic who holds that the readymade was a joke that the artworld should have declined to take seriously denies nothing the account asserts. She accepts that \emph{Fountain} was taken up and therefore is art, and holds that it is bad art, or that the practice erred in giving it the attention it did. The account cannot adjudicate this dispute, since the standards to which uptake answers are the critical practice's own. This is the limit the readymade marks: it shows what the institutional theory can explain, namely how an ordinary object becomes a work, and what it cannot, namely whether it should have.

\begin{objection}[Intention]\label{o:intent}
Duchamp may have meant \emph{Fountain} as anti-art, a gesture against the institution of art. If so, it was not presented as art, and the first condition fails.
\end{objection}

\begin{reply}
Presentation as art in the sense of Definition~2.1 of AES-15-a is offering an object as a candidate for the artworld's appreciation and criticism. An object offered as a challenge to that practice is offered to it, and invites its criticism; the challenge is one of the things critics interpret. The objection confuses the content of a work, which may be hostile to art, with its status, which depends on how it is offered and received.
\end{reply}

\section{Conclusion}

The institutional account of AES-01 implies that \emph{Fountain} is art, that its status was secured by critical uptake rather than by its submission, which was refused, and that the replicas share that status as the practice takes them up as the work. The same case marks the account's limits: it classifies the readymade without explaining its significance or vindicating the attention it has received.

\begin{thebibliography}{99}\small
\bibitem[Binkley(1977)]{Binkley1977} Binkley, T. (1977). Piece: Contra aesthetics. \emph{The Journal of Aesthetics and Art Criticism}, 35(3), 265--277.
\bibitem[Danto(1964)]{Danto1964} Danto, A.~C. (1964). The artworld. \emph{The Journal of Philosophy}, 61(19), 571--584.
\bibitem[Danto(1981)]{Danto1981} Danto, A.~C. (1981). \emph{The Transfiguration of the Commonplace: A Philosophy of Art}. Harvard University Press.
\bibitem[de Duve(1996)]{deDuve1996} de Duve, T. (1996). \emph{Kant after Duchamp}. MIT Press.
\end{thebibliography}

\end{document}
